\documentclass[11pt,a4paper]{article}
\usepackage[utf8]{inputenc}
\usepackage[T1]{fontenc}
\usepackage{amsmath,amssymb,amsthm}
\usepackage{geometry}
\geometry{margin=2.5cm}
\usepackage{hyperref}
\usepackage{booktabs}

\newtheorem{theorem}{Theorem}
\newtheorem{lemma}[theorem]{Lemma}
\newtheorem{proposition}[theorem]{Proposition}
\newtheorem{conjecture}[theorem]{Conjecture}
\newtheorem{definition}[theorem]{Definition}
\newtheorem{remark}[theorem]{Remark}

\title{A unified lower bound for the Wilf excess of numerical semigroups\\
       indexed by multiplicity and defect}
\author{(draft, preprint skeleton -- generated in-session)}
\date{\today}

\begin{document}
\maketitle

\begin{abstract}
Let $S$ be a numerical semigroup with multiplicity $m$, embedding dimension
$e$, conductor $c$ and $L = \#(S \cap [0,c-1])$. Wilf's conjecture states
$W(S) := e L - c \ge 0$. Writing $d = m - e$ for the \emph{defect}, we
propose an explicit lower bound
\[
  W(S) \;\ge\; W_{\min}(m,d) \;:=\; (m-d)\,L(d) - 2m,
  \qquad
  L(d) \;=\; \Big\lceil \tfrac{\sqrt{8d+1}-1}{2} \Big\rceil + 2,
\]
valid for every numerical semigroup $S$ with $m(S) = m$ and $d(S) = d$.
We verify the bound exhaustively for all $(m,d)$ with $m \le 24$, $d < m$
and $k^\star(S) \le 3$, by enumerating the integer points of the Kunz
polyhedron $P_m$ -- a total of $1.24 \times 10^{10}$ semigroup classes --
and we exhibit explicit sharp witnesses for every $(m,d)$ where the bound
is attained. A \emph{stabilization threshold} $m_{\min}(d)$ is observed:
for $m \ge m_{\min}(d)$ the bound is tight; for $m < m_{\min}(d)$ the
observed minimum strictly exceeds $W_{\min}(m,d)$. We record
$m_{\min}(d)$ for $d \le 13$ and give a conjectural growth rate.
\end{abstract}

\section{Introduction}

Wilf's conjecture \cite{Wilf1978}, open since 1978, asks whether
every numerical semigroup $S$ satisfies $e L \ge c$, equivalently
$W(S) \ge 0$. The conjecture is known in several regimes: $e \le 3$
(Froberg-Gottlieb-Haggkvist), $e \ge m/2$ (Sammartano, Eliahou),
and small genus (Bras-Amoros, Fromentin-Hivert). The present paper
does not attempt a proof. Instead, we investigate the \emph{shape} of
the quantity $W(S)$ as a function of the two coarsest invariants,
multiplicity $m$ and defect $d := m - e$, and we conjecture an
explicit two-parameter lower bound which we verify computationally
over a large window of the Kunz polyhedron.

\paragraph{Motivation.}
Writing $W$ as a function of $(m,d)$ rather than $(e,g)$ has two
advantages: (i) it matches the natural axes of the Kunz polyhedron
$P_m$ and therefore of any enumeration algorithm organized by
multiplicity; (ii) it exposes a striking regularity in the minimum,
namely that $\min_{S : (m,d)} W(S)$ eventually depends on $d$ alone
through a simple $L(d)$ factor, independent of $m$ beyond a threshold.

\section{Definitions and notation}

\begin{definition}
For a numerical semigroup $S \subsetneq \mathbb{N}$ with $0 \in S$,
\begin{itemize}
  \item $m = m(S) = \min (S \setminus \{0\})$ is the \emph{multiplicity};
  \item $e = e(S)$ is the number of minimal generators (\emph{embedding dimension});
  \item $F = F(S) = \max(\mathbb{Z} \setminus S)$ is the \emph{Frobenius number}, $c = F + 1$;
  \item $L = L(S) = \#(S \cap [0,c-1])$;
  \item $d = d(S) := m - e$ is the \emph{defect};
  \item $W(S) := e L - c$ is the \emph{Wilf excess}.
\end{itemize}
\end{definition}

The Kunz coordinates $(k_1,\ldots,k_{m-1})$ of $S$ are defined by
$w_i = k_i m + i$, where $w_i = \min\{s \in S : s \equiv i \pmod{m}\}$
is the Apery set. We write $k^\star = \max_i k_i$; this controls the
conductor via $c \le k^\star \cdot m + (m-1)$.

\section{The unified conjecture}

\begin{conjecture}[Unified lower bound]\label{conj:unified}
For every numerical semigroup $S$ with $m(S) = m \ge 2$ and $d(S) = d \in
\{0,1,\ldots,m-1\}$,
\[
   W(S) \;\ge\; W_{\min}(m,d) \;:=\; (m-d)\,L(d) - 2m,
\]
where
\[
  L(d) \;=\;
  \begin{cases}
    2 & \text{if } d = 0,\\[2pt]
    \Big\lceil \dfrac{\sqrt{8d+1} - 1}{2} \Big\rceil + 2 & \text{if } d \ge 1.
  \end{cases}
\]
\end{conjecture}

\begin{remark}
$L(d)$ is the unique increasing step function such that
$L(d) = \ell + 2$ on the triangular interval
$d \in \big[\tfrac{\ell(\ell-1)}{2} + 1,\ \tfrac{\ell(\ell+1)}{2}\big]$.
Its appearance is not arbitrary: in each step, the bound corresponds
to Apery sets concentrated at $k^\star \in \{\ell+1,\ell+2\}$, and the
triangular numbers mark the exact defects at which $k^\star$ must
increase.
\end{remark}

\begin{remark}
Conjecture~\ref{conj:unified} implies the classical Wilf bound
$W(S) \ge 0$ whenever $(m-d) L(d) \ge 2m$, which holds
unconditionally for $d \ge 0$ once $L(d) \ge 2$ and $m \ge d$.
Hence Conjecture~\ref{conj:unified} refines Wilf's conjecture in the
sense that its truth on $(m,d)$ with $d < m$ entails
$W(S) \ge 0$ on the same window.
\end{remark}

\section{Stabilization threshold}

For each defect $d$ we define
\[
  m_{\min}(d) \;:=\; \min\big\{\,m \ge d+1 :
    \min_{S : m(S)=m,\, d(S)=d} W(S) = W_{\min}(m,d)\,\big\}.
\]
Our experiments (Section~\ref{sec:num}) give Table~\ref{tab:mmin}.

\begin{table}[h]
\centering
\begin{tabular}{c|cccccccccccccc}
\toprule
$d$         & 0 & 1 & 2 & 3 & 4 & 5 & 6 & 7 & 8 & 9 & 10 & 11 & 12 & 13\\
\midrule
$m_{\min}(d)$ & 2 & 3 & 4 & 6 & 8 & 10 & 13 & 15 & 18 & 21 & 23 & $\le 24$ & $\le 24$ & $\le 24$\\
\bottomrule
\end{tabular}
\caption{Observed stabilization thresholds. For $d \le 10$ the values
are confirmed to be sharp (i.e.\ the bound fails to be attained at
$m = m_{\min}(d) - 1$). For $d \in \{11,12,13\}$ we only know the
bound becomes sharp at $m = 24$; confirming it as a true threshold
requires a complete enumeration at $m \in \{22, 23\}$.}
\label{tab:mmin}
\end{table}

\begin{conjecture}[Growth of the threshold]\label{conj:mmin}
$m_{\min}(d) = 2d + O(\sqrt d)$; more precisely,
$m_{\min}(d) - 2d$ is monotone and grows like $\sqrt{2d}$.
\end{conjecture}

\section{Sharp witnesses}

For each $(m,d)$ with $m \in \{20,22,23,24\}$ and $d < m$
where the bound is attained ($68$ cases total) we extracted a
concrete witness $S(m,d)$ directly from the argmin of the enumeration.
These witnesses are catalogued in
\texttt{results/sharp\_witnesses.json}; Table~\ref{tab:witnesses}
reproduces a representative subset.

\begin{table}[h]
\centering
\small
\begin{tabular}{cccccl}
\toprule
$m$ & $d$ & $e$ & $L$ & $W$ & minimal generators\\
\midrule
20 & 2 & 18 & 4 & 32 & $\langle 20,21,22,45,\ldots \rangle$\\
20 & 9 & 11 & 6 & 26 & $\langle 20,21,22,26,29,45,53,\ldots \rangle$\\
22 & 5 & 17 & 5 & 41 & $\langle 22,23,24,27,52,\ldots \rangle$\\
23 &10 & 13 & 6 & 32 & $\langle 23,24,27,32,34,49,\ldots \rangle$\\
24 & 5 & 19 & 5 & 47 & $\langle 24,25,26,29,56,\ldots \rangle$\\
24 &13 & 11 & 7 & 29 & $\langle 24,25,27,28,33,35,64,\ldots \rangle$\\
\bottomrule
\end{tabular}
\caption{Representative sharp witnesses. Each row satisfies
$W = W_{\min}(m,d)$ exactly.}
\label{tab:witnesses}
\end{table}

\begin{proposition}[empirical, Section~\ref{sec:num}]
For every $(m,d)$ listed in Table~\ref{tab:mmin} with
$m \ge m_{\min}(d)$, at least one sharp witness exists.
\end{proposition}

The existence of explicit witnesses rules out the possibility that
$W_{\min}(m,d)$ is a strict lower bound only by accident of the
enumeration window, and is essential for any future attempt at a
closed-form construction.

\section{Numerical verification}
\label{sec:num}

\paragraph{Enumeration.}
We enumerate the integer points of the Kunz polyhedron
$P_m \subset \mathbb{R}^{m-1}$ by backtracking over
$k_1, k_2, \ldots, k_{m-1}$, forward-checking the ``no-carry''
inequalities $k_r \le k_a + k_b$ ($a+b = r$) and verifying the
``carry'' inequalities $k_r \le k_a + k_b + 1$ ($a+b = r+m$) at
leaves. For each leaf we compute $(d, L, c, W)$ from the
Apery set in $O(m^2)$ time.

\paragraph{Scale.}
The enumeration is implemented in C (\texttt{kunz\_core.c}) with
an OpenMP variant parallelizing the outer
$(k_1, k_2)$ loop. On a 4-core Xeon the parallel kernel gives
a $\approx 2\times$ speedup over the serial kernel. For
$K_{\max} = 3$ (i.e.\ $k^\star \le 3$), we completed the
enumeration for every $m \in \{20, 22, 23, 24\}$ and obtained
$1.24 \times 10^{10}$ valid semigroup classes across
$70$ distinct $(m,d)$ slices, with \emph{zero} violations of
Conjecture~\ref{conj:unified}.

\paragraph{Cross-checks.}
(a) Internal: an independent BFS+Apery implementation agreed with
the Kunz enumerator on $1000/1000$ randomly sampled semigroups.
(b) Targeted: a dedicated counter-example hunt in the regime
$d > 2m/3$ (where the gap to $W_{\min}(m,d)$ is smallest) found
no violation. (c) Bound consistency: the unified formula is
tight for $d \in \{0,1,2\}$ at $m \in \{2,3,4\}$ respectively,
i.e.\ at the smallest $m$ where each defect is realizable.

\section{Open problems}

\begin{enumerate}
  \item Prove Conjecture~\ref{conj:unified}. A natural strategy is
  to fix $d$ and argue that the Kunz polytope $P_{m,d}$ contains a
  rational vertex attaining $W = W_{\min}(m,d)$ for $m$ large enough,
  then to bound $W$ on all other vertices.
  \item Determine $m_{\min}(d)$ exactly as a closed-form expression
  in $d$ (Conjecture~\ref{conj:mmin}).
  \item Extend the verification beyond $k^\star = 3$. Results at
  $(m,K_{\max}) = (20,4)$ match those at $(20,3)$ for every
  $(m,d)$, suggesting that the minimum is already captured by
  $k^\star \le 3$; a proof of this \emph{$k^\star$-saturation} would
  reduce the enumeration complexity substantially.
  \item Characterize, for each $(m,d)$, the full set of sharp
  witnesses. Is it a union of one-parameter families?
\end{enumerate}

\bibliographystyle{alpha}
\begin{thebibliography}{9}
\bibitem{Wilf1978} H.S. Wilf, \emph{A circle-of-lights algorithm for
the ``money-changing problem''}, Amer. Math. Monthly 85 (1978),
562--565.
\end{thebibliography}

\end{document}
